\scanpage{38}
18)\quad Les foncteurs entre cat. triangulés \deleted{\unsure{hom}}
commutent aux foncteurs translation et transforment triangles distingués
en \unsure{idem} ....

\[
D(C),\quad D^-(C)\qquad
\begin{gathered}\text{importants = sans qu'il y}\\\text{ait des projectifs.}\end{gathered}
\]
\[
D^b(C)
\]
\diagram{\begin{tikzcd}[column sep=large]
Lf^*:D(Y)\arrow[r]&D(X)\\
D^+(Y)\arrow[r,dash]&D^+(X)\\
D^-(Y)\arrow[r,dash]&D^-(X)
\end{tikzcd}}
quand on a $f:X\mathbin{-}Y$ morphisme \uuline{plat}
de topos \uline{annelés}

\[
D^+(Y)\mathbin{-}D^+(X)
\]
définition directe via résolutions injectives

\uline{ou} définition indirecte par ``passage au quotient''

[Pour calculer $Lf^*(K)$, pas la peine de remplacer au préalable $K$
par une résolution injective.]

\uuline{N.B.} Si $f$ n'est \uline{pas plat}, on \deleted{\unsure{f}} ne peut
en général définir raisonnablement
\[
f^*:D^+(Y)\longrightarrow D^+(X)
\]
(de sorte à avoir des bonnes propriétés).
\margin{Rien n'empêche de construire un dérivé droit, $Rf^*$, mais il n'a pas
le moindre intérêt.}
Mais dans la suite de ces idées, on explique comment construire raisonnablement
\[
Lf^*:D^-(Y)\mathbin{-}D^-(X),
\]
\inserted{``dérivé gauche'' de $f^*$,}
en utilisant (\unsure{non par des rés. projectives mais}) des résolutions plates.
Ce n'est que si $f$ est de ``Tor-dim. finie'' que cela définit
\[
D^+(Y)\longrightarrow D^+(X)\ ;\quad D^b(Y)\mathbin{-}D^b(X).
\]
\mspage{37}
